\chapter{Introduction and Project Vision}
\label{chap:intro}
%=====================================================================

\section{Project Overview}
Aegis is an egocentric spatio-temporal vision system that helps people with
cognitive impairment recall where objects are and what they did earlier in the
day, inside indoor environments such as the home. A wearable first-person
camera continuously observes the user's surroundings. Aegis converts these
observations into a persistent spatio-temporal memory: a structured record of
objects, their locations relative to stable parts of the environment, the
layout of the space, and the activities that took place, all indexed by time.

To remain usable over weeks and months, the memory is \emph{multi-resolution}.
Recent events are stored in fine spatial and contextual detail, while older
events are progressively condensed into compact, higher-level temporal
summaries. The user, or an authorised caregiver, asks questions in natural
language by speech or text (``Where did I leave my glasses?'', ``Did I take my
tablets this morning?''). Aegis grounds the answer in the stored memory and
shows the supporting evidence.

Aegis is an assistive tool. It does not replace human care and is not a
certified medical device (see Section~\ref{sec:ethics}).

\section{Problem Statement}
People with subjective cognitive decline, mild cognitive impairment, or
dementia routinely lose track of where they put objects and whether they have
done daily tasks. To the best of our knowledge, none of the systems reviewed in
Chapter~\ref{chap:research} maintains a long-term, spatially grounded, queryable record of
first-person experience that stays computationally manageable as history
accumulates and that also gives caregivers verifiable access without constant
supervision. As a result, users search manually, repeat searches, or ask a
caregiver who often was not present and must reconstruct events from memory.

\noindent\textbf{Formal problem statement.} Given a continuous, unconstrained
egocentric video stream from an indoor environment, construct and maintain a
persistent spatio-temporal memory that (a) preserves recent spatial detail,
(b) compresses older history without destroying information users are likely
to need, and (c) answers natural-language spatial and temporal queries with
evidence-grounded responses at low latency, for both a primary user with
cognitive impairment and a delegated caregiver.

\section{Problem Context and Motivation}
The problem unfolds as a repeating loop (Figure~\ref{fig:recall-loop}): an
object is misplaced or an activity is forgotten and no record exists; the
person searches manually, in a disorganised and often repeated way; the person
asks a caregiver, who reconstructs the answer from memory, often without having
been present; the caregiver searches as well, duplicating effort, and
resolution comes by luck or escalation.

\aegisfig{fig-1.1}
  {The current ``no record'' recall loop, and where Aegis breaks it.}
  {fig:recall-loop}

Each step of the loop carries a cost: manual searching consumes time and
effort; asking a caregiver interrupts them and may still fail, because the
caregiver often was not present; duplicated searching multiplies the work for
both parties. Everyday memory lapses of this kind are common among older adults
\cite{mogle2023,marchant2011}, and dementia places a disproportionate burden on
family and other unpaid caregivers \cite{kasper2015}. Aegis breaks the loop at
its first step: by keeping a persistent, queryable record, the manual search
and the caregiver's reconstruction are replaced by a direct answer with
evidence.

\subsection*{Scale and significance of the problem}
\begin{itemize}
\item \textbf{Ageing population.} The World Health Organization projects about
2.1 billion people aged 60 and over worldwide by 2050 \cite{who2024}. In the
United States the share of seniors is projected to rise from 16\% to 22\%
\cite{statista2022}.
\item \textbf{Memory decline is widespread.} Subjective cognitive decline, mild
cognitive impairment, and dementia affect a substantial share of older adults
\cite{alz2023}. Daily memory lapses, including losing objects, are common and can be
disruptive \cite{mogle2023,marchant2011}.
\item \textbf{Ageing in place.} Independent living is widely preferred, and some
people with early-stage Alzheimer's disease live alone and voice concerns about
coping without support \cite{harris2006}.
\item \textbf{Caregiver burden.} Informal caregiving imposes substantial economic costs, with opportunity-cost estimates for elder care in the United States around \$522 billion per year \cite{chari2015}.
\end{itemize}

Several of these figures come from United States sources. The motivation is
therefore to give the person a way to recall independently, reducing how often
they must depend on someone else, while giving caregivers evidence-based
answers instead of guesses. 

\section{Project Vision}
Aegis aims to be a quiet, dependable ``external episodic memory'': a system
where a person with cognitive impairment can ask a plain question and receive a
short, correct, evidence-backed answer without learning an interface, and one
that a caregiver can trust for oversight without watching the person
continuously. The long-term vision is a system that scales from days to months
of use on modest hardware, degrades gracefully (saying ``I'm not sure'' rather
than guessing), and respects the dignity and privacy of the person it serves.

\section{Project Objectives}
Table~\ref{tab:objectives} lists the project objectives with their proposed
measures. The targets below are deliberately lenient minimum-viable thresholds for a first working prototype. They are not production safety specifications and may be raised in later iterations.

\begin{table}[H]
\centering\footnotesize
\caption{Project objectives, measures and status at the end of FYP-1.}
\label{tab:objectives}
\begin{tabularx}{\textwidth}{@{}>{\bfseries}l L{3.2cm} Y Y@{}}
\toprule
\rowcolor{hdrfill}\textbf{ID} & \textbf{Objective} & \textbf{Measure (proposed target)} & \textbf{Status at FYP-1 end}\\
\midrule
O1 & Continuously capture and ingest egocentric video into the processing pipeline. & Uninterrupted ingestion for at least 1 hour without crash or stall. & \textbf{Partial}: live capture with reconnect, gaps, pause and resume implemented and tested (\path{ingestion/live_pipeline.py}, 64 ingestion tests); 1-hour soak not run (P2.1).\\
O2 & Build a spatio-temporal memory of room regions, object locations, layout and activities. & Room segmentation accuracy $\geq$ 0.90; room-grounded correct-room rate $\geq$ 0.90 (Chapter~\ref{chap:testing}). & \textbf{Partial}: rooms, anchors, relations, state runs and activities implemented; room metrics \textbf{Unmeasured} (P2.3).\\
O3 & Progressively compress older history while retaining recent detail. & Sub-linear memory growth; accuracy drop on compressed history $\leq$ 15 points. & \textbf{Partial}: compaction implemented and verified; 5.7--7.2$\times$ on short clips; accuracy drop after compaction \textbf{Unmeasured}.\\
O4 & Answer natural-language spatial and temporal queries with grounded responses. & Spatial accuracy $\geq$ 70\% and temporal accuracy $\geq$ 60\%; every answer cites at least one evidence item. & \textbf{Partial}: spatial queries answered with evidence; temporal \textbf{Planned} (the engine declines); accuracy \textbf{Unmeasured} (\texttt{demo05} not annotated).\\
O5 & Support speech and text queries through an accessible interface. & A first-time user completes a query within 60 s without help. & \textbf{Partial}: on-device speech and text in both apps; usability study not run.\\
O6 & Provide authorised, consent-based caregiver access. & Unauthorised attempts rejected in 100\% of test cases. & \textbf{Planned}: Caregiver App exists; no authentication yet (P1.7).\\
O7 & Evaluate system efficacy against structured ground truth. & Evaluation dataset selected and documented. & \textbf{Partial}: evaluation harness built; EgoLife used for the multi-day study; Ego4D NO-GO (0/100); held-out \texttt{demo05} NOT\_READY.\\
\bottomrule
\end{tabularx}
\end{table}

\section{Project Scope}
\textbf{In scope:} indoor, home-like environments; egocentric (first-person)
visual input as the sole sensing modality; object-location and activity
tracking in a persistent spatio-temporal memory; multi-resolution memory
abstraction; natural-language query answering over stored memory with speech
and text input; a user-facing assistive interface and a delegated caregiver
application; room identity calibration for one home; evaluation of spatial and temporal retrieval accuracy against structured ground truth. Evaluation now uses a public seven-day egocentric dataset (EgoLife, participant A1\_JAKE) for multi-day persistence, alongside a frozen held-out care-style clip (\texttt{demo05}); Ego4D was assessed and is not pursued in FYP-2 unless access, data and an official scorer become available (Section~\ref{sec:egolife}). The per-home registry preserves room identity for a single home; it does not provide multi-home deployment support.

\textbf{Out of scope:} outdoor navigation and wayfinding; medical diagnosis,
disease staging, or treatment recommendations; emergency detection (falls,
wandering) and real-time alerts; non-visual sensing (biosensors, smart-home
sensors, RFID or BLE tags); certification as a medical device and clinical
trials with patients; multi-user or multi-home deployments.

\begin{table}[H]
\centering\small
\caption{Scope split across FYP-1 and FYP-2.}
\label{tab:scope}
\begin{tabularx}{\textwidth}{@{}L{3.6cm} Y Y@{}}
\toprule
\rowcolor{hdrfill}\textbf{Scope item} & \textbf{FYP-1} & \textbf{FYP-2}\\
\midrule
Capture and ingestion (M1) & Iterations 1--2: batch, live capture, keyframes, gaps, segments & Hardening: 1-hour soak, disconnect recovery measured\\
Spatio-temporal memory and three-day retention (M2) & Iterations 1--2: rooms, relations, state runs, home store, compaction & Room-identity fix, room metrics, long-term persistence study\\
Query engine (M3) & Iteration 2: closed-vocabulary location queries, ranking, abstention, on-device speech & Iteration 3: temporal and activity answers, calibrated confidence, evidence references\\
Caregiver access, security, performance (M4) & Iteration 2: Caregiver App (unauthenticated) & Iteration 4: authentication, grants, audit, encryption, latency and usability validation\\
\bottomrule
\end{tabularx}
\end{table}

Figure~\ref{fig:context} fixes the system boundary of the FYP-1 prototype.
The primary user and caregiver are outside the boundary and reach Aegis through
two Android applications. The wearer's phone acts only as a camera, streaming
MJPEG over Wi-Fi. Heavy computation runs on Modal, and model weights come from
the Hugging Face Hub. Speech is transcribed on the user's phone by the Android
speech service, so no audio reaches Aegis. The indoor environment is external
and observed only through the camera; the note ``includes bystanders'' makes the
bystander-privacy concern of Section~\ref{sec:ethics} visible at the boundary.
Authentication and consent services are drawn dashed because they are planned
for FYP-2. The data flows shown here seed the interface tables of
Chapter~\ref{chap:design}.

\aegisfig{ch1-system-context}
  {System context of Aegis at the end of FYP-1: external entities and the data
   crossing the system boundary. Dotted: physical influence; dashed: planned.}
  {fig:context}

\section{Project Modules}
\label{sec:modules}
The system is organised into four modules, aligned with the team's division of
work (Table~\ref{tab:modules} and Figure~\ref{fig:modules}).

\begin{table}[H]
\centering\small
\caption{Module decomposition.}
\label{tab:modules}
\begin{tabularx}{\textwidth}{@{}l L{2.6cm} Y l l@{}}
\toprule
\rowcolor{hdrfill}\textbf{ID} & \textbf{Name} & \textbf{Responsibility} & \textbf{FRs} & \textbf{Lead}\\
\midrule
M-A & Capture and Ingestion & Live phone capture, keyframe selection, gap/pause records, crash resume, promotion into bounded segments, batch video normalisation. & FR1, FR2 & Shahmeer\\
M-B & Spatio-Temporal Memory & Gravity frame, 3D lifting, relations, state runs, events, rooms, per-home store, cross-session matching, activities, evidence, compaction. & FR3, FR4 & Faizan\\
M-C & Query Engine & Closed-vocabulary NLQ (\path{src/aegis/nlq/}), ranked candidate retrieval, abstention, answer wording; \path{/api/v1/ask}. & FR5, FR6 & Ali Sher\\
M-D & Interfaces and Caregiver Access & Patient and Caregiver apps, API routes, deployment; authentication planned. & FR7, FR8 & Shahmeer\\
\bottomrule
\end{tabularx}
\end{table}

\aegisfig{ch1-modules}
  {Module overview: four modules, their functional requirements and leads,
   and the shared memory store. \tbd{Leads to be confirmed by the team.}}
  {fig:modules}

\textbf{Module responsibilities.} M-A turns the phone's MJPEG stream into a
disciplined feed of timestamped keyframes and owns ingestion reliability. M-B
owns the memory: it converts reconstructed geometry and detections into rooms,
anchors and objects, records each object's stable states and the events between
them, merges sessions into a persistent per-home store and compacts history
older than three days into verified summaries. M-C owns understanding: it maps
a question to a known object using a closed vocabulary, retrieves and ranks
candidates from memory, and composes a short answer from a template, or
declines when evidence is missing. No language model writes answer text. M-D
owns access: the patient and caregiver applications and the HTTP API. Leads are
responsibility assignments; Iteration~2 work crossed module boundaries, and
Chapter~\ref{chap:pm} records who built what.

\section{Stakeholders and User Classes}
\begin{table}[H]
\centering\footnotesize
\caption{Stakeholders and user classes.}
\label{tab:stakeholders}
\begin{tabularx}{\textwidth}{@{}L{2.2cm} Y Y Y@{}}
\toprule
\rowcolor{hdrfill}\textbf{Class} & \textbf{Description} & \textbf{Needs} &
\textbf{Characteristics and design implications}\\
\midrule
Primary user & Person with cognitive impairment (SCD, MCI, or early dementia)
living at home. & Independent recall of object locations and daily activities. &
May have reduced working memory, low technical confidence, possible hearing or
vision limits. Needs short, simple, low-load answers and speech input.\\
Secondary user: caregiver & Family member or professional carer. & Verifiable
oversight without constant supervision. & Needs history queries, object-state
verification, and activity summaries; must act only with the primary user's
consent.\\
Supervisor and evaluators & Academic stakeholders. & Evidence of rigour,
reproducibility, and responsible tool use. & Need traceability from
requirements to tests.\\
Bystanders & Other people visible in the camera view (visitors, household
members). & Privacy. & Not users, but affected; drives the privacy
requirements.\\
\bottomrule
\end{tabularx}
\end{table}

\section{Success Criteria}
\label{sec:success}
Success at the end of FYP-2 will be judged by the following measurable
criteria. These thresholds are minimum-viable prototype targets, deliberately
lenient for FYP-stage validation and not production safety specifications. Numeric thresholds are proposed and will be finalised with the
supervisor once a benchmark is fixed.

\begin{table}[H]
\centering\footnotesize
\caption{Success criteria for FYP-2 and the evidence available at the end of FYP-1.}
\label{tab:success}
\begin{tabularx}{\textwidth}{@{}>{\bfseries}l Y L{2.9cm} Y@{}}
\toprule
\rowcolor{hdrfill}\textbf{ID} & \textbf{Criterion} & \textbf{Threshold} & \textbf{FYP-1 evidence}\\
\midrule
SC1 & Spatial retrieval accuracy (``where is X?'') & $\geq$ 70\% & \textbf{Unmeasured}: held-out \texttt{demo05} annotations missing (NOT\_READY)\\
SC2 & Temporal retrieval accuracy (``when / what did I do?'') & $\geq$ 60\% & \textbf{Unmeasured}: temporal intent not implemented\\
SC3 & Accuracy retained after compaction & Drop $\leq$ 15 points & \textbf{Partial}: 100\% agreement on auto-generated probes (demo03/04); human queries unmeasured\\
SC4 & End-to-end query latency & $\leq$ 5\,s at p95 & \textbf{Unmeasured} (P2.8)\\
SC5 & Evidence grounding & 100\% of non-abstaining answers & \textbf{Partial}: \path{/ask} and \path{/query} return evidence; coverage not audited (P1.4)\\
SC6 & Safe failure: explicit ``not sure'' instead of a fabricated answer & Verified on unanswerable queries & \textbf{Partial}: abstention unit tests (\texttt{test\_nlq.py}); EgoLife keys/wallet abstention probes; calibrated threshold planned (P1.2)\\
SC7 & Caregiver access control & No unauthorised access & \textbf{Not met}: no authentication yet (P1.7)\\
SC8 & Usability & 5 participants unaided & \textbf{Unmeasured}: needs ethics approval\\
SC9 & Room-level spatial memory & Segmentation $\geq$ 0.90, room-change F1 $\geq$ 0.70, correct-room $\geq$ 0.90 & \textbf{Unmeasured} (P2.3); identity instability observed on EgoLife\\
\bottomrule
\end{tabularx}
\end{table}

\section{Constraints and Assumptions}
\label{sec:constraints}
\textbf{Constraints.}
\begin{enumerate}[label=\textbf{C\arabic*.},leftmargin=3em]
\item \emph{Academic timeline:} FYP-1 (design, early implementation) and FYP-2
(completion, evaluation) within the 2023--2027 session.
\item \emph{Compute and budget:} development on student hardware and
three student laptops: two Lenovo T480 systems running Arch Linux and one HP laptop of approximately similar specifications running Debian/Ubuntu (laptop CPU, 8~GB RAM, about 100~GB free SSD each). The capture device is a mobile phone with an RGB-only camera. Model inference uses Modal L4 GPUs; no model is trained or fine-tuned. One five-minute segment costs about 0.086 GPU-hours (about US\$0.069: VGGT 200\,s, detection 108\,s); the 74-segment EgoLife study cost US\$8.71 metered. A CPU fallback (\texttt{AEGIS\_GPU=cpu}) runs the same stages without a GPU at much higher latency (VGGT 831\,s for 79 frames).
\item \emph{Clinical scope:} clinical trials with patients are out of scope
(Section~1.6); validation relies on healthy volunteers and scripted scenarios.

\item \emph{Data:} EgoLife (A1\_JAKE, seven days, 12:00--13:00) is used for the multi-day study; a held-out care-style clip (\texttt{demo05}) is frozen for accuracy scoring; Ego4D was assessed as NO-GO (0/100 readiness, about 10.8 weeks of prerequisites).
\item \emph{Home identity:} the room registry persists calibration for one
home identified by \texttt{home\_id}; it does not constitute multi-home service
or support. 
\item \emph{Team size:} three members, each owning separate modules with shared
milestones.
\end{enumerate}

\textbf{Assumptions.}
\begin{enumerate}[label=\textbf{A\arabic*.},leftmargin=3em]
\item The user carries a mobile phone as the first-person RGB capture device during the test session; setting up a dedicated camera is out of scope.
\item The camera view approximates the user's attention and shows the objects
they handle or place.
\item Indoor lighting is adequate for object recognition most of the time.
\item Network connectivity is available for any cloud-hosted model calls, or a
local fallback is provided. The live prototype uses local storage and a 2 FPS sampling target.
\item A caregiver, where involved, is nominated by the primary user (or legal
guardian) and authorised before access begins.
\item Objects moved while the user is not wearing the camera cannot be tracked;
the system reports this limitation rather than guess.
\item VGGT-$\Omega$ coordinates are used without manual axis permutation or
scale factor. Room footprints use floor-plane geometry, but monocular
reconstruction and inferred room boundaries are approximate and do not
establish metric accuracy.
\item Only an RGB camera is assumed: no depth sensor or external pose tracking.
This follows the RGB-only wearable setting of R4DSG \cite{r4dsg2026} rather
than the RGB-D-plus-pose input of DAAAM \cite{daaam2026}.
\item VGGT-$\Omega$ weights are taken from \path{facebook/VGGT-Omega} when access is granted, otherwise from the public mirror \path{wincentIsMe/VGGT-Omega} (same weights and licence).
\end{enumerate}

\textbf{Feasibility.} The building blocks of Aegis exist separately in recent
work: MemPal shows that voice-based object retrieval from a wearable camera
helps older adults in their own homes \cite{mempal2025}; R4DSG shows that
persistent, anchor-relative object memory can be built from RGB-only wearable
video \cite{r4dsg2026}; and DAAAM shows that a tool-calling agent can answer
spatio-temporal questions over a hierarchical memory \cite{daaam2026}. The
project's novelty is their integration with progressive compaction and
consent-based caregiver access, within a deliberately limited scope: indoor,
single user, camera-only. Each module has one owner (Section~\ref{sec:responsibilities}), and if
tracking proves unreliable in cluttered scenes the evaluated environment will
be restricted to fewer anchors and objects (Section~\ref{sec:deviations}).

\section{Ethical, Legal and DEI Considerations}
\label{sec:ethics}
\textbf{Societal, economic, environmental.} Aegis supports independence and
dignity for the person and reduces caregiver burden. By lowering the need for
constant supervision it may reduce long-term care costs. Lightweight
processing and keyframe selection limit compute and energy use.

\textbf{Ethics.} Aegis is designed to \emph{assist, not replace}: it supports
human caregiving and never substitutes for it. Safe failure is a first-class
requirement: when evidence is weak or conflicting, Aegis says it is unsure
rather than inventing an answer, because a wrong confident answer (for example,
a wrongly stated medication event) can cause harm. The caregiver portal exists
for verifiable help the user has agreed to, \emph{not} for covert monitoring;
access is consent-based, scoped, and auditable. The primary user can see what
is stored and who has access. 

\textbf{Privacy and data protection.} Continuous first-person video captures
highly sensitive material, including bystanders and private spaces. Design
commitments: process and store data locally where feasible; store derived
structured memory rather than raw video where possible; apply retention limits
through compaction; encrypt data at rest and in transit; provide a means to
pause capture and delete history; and avoid capturing or retaining identifiable
bystander imagery beyond what is needed.

\textbf{Legal.} Aegis is not a certified medical device; outputs are advisory,
not final decisions. Informed consent is required from the primary user (or an
authorised representative where capacity is limited) and, as far as
practicable, from other household members whose images may be captured.

Test videos used by the team include only consenting participants; bystanders are excluded or their identifiable imagery is removed before analysis.

\textbf{Data used in FYP-1.} Team recordings include only consenting team
members. EgoLife is a public dataset of consenting volunteers and is used for
evaluation only; its frames are processed on Modal and not redistributed.
Object-only evidence crops (at most 320\,px) are kept so answers can show a
photo; this exceeds the 24-hour raw-frame rule in NFR-S4 and is recorded as
supervisor decision \#9 \tbd{(supervisor decision pending)}; it can be disabled
with \texttt{memory.evidence.enabled}.

\textbf{Current security status.} The API has no authentication, allows any
origin (CORS \texttt{*}) and both apps permit cleartext HTTP on a LAN for
development. This is acceptable only for a single-user prototype with test
data and is the main FYP-2 security work (Section~\ref{sec:security}).

\textbf{DEI and accessibility.} Speech and text input serve users with
different abilities and literacy levels. The interface should use large text,
high contrast, and simple wording suitable for older adults. Users with speech
or hearing difficulties must be able to use text and visual output alone.
Language support is English-only; supporting Urdu is a candidate
extension.  Evaluation should include participants of different
ages and, where possible, both genders, so accuracy is not tuned to one
group's environment or appearance.

%=====================================================================
